% Generated by src/equivalent_rights_proposals.jl -- do not edit by hand.
\begin{table}[htbp]
\centering
\small
\caption{Club members left below their welfare under the proposal, by variant}
\begin{tabular}{lcccccccc}
  \toprule
  & \multicolumn{2}{c}{\textbf{Wolfram et al.}} (of 47) & \multicolumn{2}{c}{\textbf{Banerjee, Duflo \& Greenstone}} (of 179) & \multicolumn{2}{c}{\textbf{Equal Right}} (of 173) & \multicolumn{2}{c}{\textbf{Equal Right, 5\%/yr}} (of 173) \\
  \cmidrule(lr){2-3} \cmidrule(lr){4-5} \cmidrule(lr){6-7} \cmidrule(lr){8-9}
  \textbf{Variant} & \textbf{isolated} & \textbf{joint} & \textbf{isolated} & \textbf{joint} & \textbf{isolated} & \textbf{joint} & \textbf{isolated} & \textbf{joint} \\
  \midrule
  \multicolumn{9}{l}{\textit{Members losing on equally-distributed-equivalent consumption}} \\
  \quad Variant 1 & 24 & 2 & 38 & 37 & 30 & 8 & 28 & 29 \\
  \quad Variant 2 & 6 & 2 & 50 & 36 & 0 & 0 & 1 & 13 \\
  \quad Variant 3 & 1 & 2 & 41 & 26 & 6 & 19 & 18 & 29 \\
  \quad Variant 4 & 6 & 9 & 50 & 69 & 0 & 0 & 1 & 53 \\
  \midrule
  \multicolumn{9}{l}{\textit{Members losing on mean consumption per capita}} \\
  \quad Variant 1 & 29 & 3 & 37 & 4 & 30 & 6 & 36 & 5 \\
  \quad Variant 2 & 2 & 0 & 10 & 0 & 22 & 0 & 0 & 0 \\
  \quad Variant 3 & 0 & 1 & 36 & 14 & 0 & 11 & 0 & 15 \\
  \quad Variant 4 & 2 & 0 & 10 & 0 & 22 & 0 & 0 & 0 \\
  \bottomrule
\end{tabular}
\label{tab:equiv_rights_losers_er}
\\[4pt]
{\footnotesize Note: a member counts as losing when its NPV of consumption over 2030--2100 falls short of what it obtains under the proposal itself. Variant 1 makes every member indifferent by construction, so any count there is numerical noise. Variants 2 to 4 return the same surplus of rights under different sharing rules: a split chosen to minimise the largest relative shortfall (2), one proportional to population times marginal utility (3), and uniform scaling (4).}
\end{table}
